\section{Modern Video Coding}




\begin{questionbox}{Domanda 89}
(ModernVC) What is the specific scope of video compression standards like H.266/\textbf{VVC}?
\end{questionbox}


Video coding standards define bitstream syntax and normative decoder behavior. They do not prescribe the full encoder implementation. \textbf{VVC} targets much higher compression efficiency than \textbf{HEVC}, especially for high-resolution, HDR, 360-degree and immersive video, at the cost of much higher complexity.



\begin{questionbox}{Domanda 90}
(ModernVC) Explain the advantage of the “Coding Tree Unit” (CTU) structure introduced in \textbf{HEVC}/\textbf{VVC}.
\end{questionbox}


Coding Tree Units replace fixed macroblocks with flexible recursive partitioning. Large homogeneous regions can use large blocks, while detailed regions can be split into smaller blocks. This improves \textbf{\textbf{rate}-\textbf{distortion}} efficiency.

\textbf{HEVC} uses quadtree partitioning. \textbf{VVC} extends this with multi-type trees, including binary and ternary splits.



\begin{questionbox}{Domanda 91}
(ModernVC) What are the roles of VCL and NAL in modern video standards?
\end{questionbox}


\begin{itemize}
  \item \textbf{VCL (Video Coding Layer):} contains compressed video data such as slices, prediction, residuals and coding syntax.
  \item \textbf{NAL (Network Abstraction Layer):} wraps VCL and non-VCL data into NAL units for transport/storage. It separates codec syntax from packetization.
\end{itemize}

Examples of non-VCL NAL units: SPS, PPS and SEI.



\begin{questionbox}{Domanda 92}
(ModernVC) What is the main purpose of the \textbf{CABAC} \textbf{entropy} coder in modern standards?
\end{questionbox}


\textbf{CABAC} means \textbf{Context-Adaptive Binary Arithmetic Coding}. It binarizes syntax elements and arithmetic-codes them with probabilities adapted to local context. It improves compression efficiency over simpler VLC/\textbf{CAVLC} while remaining \textbf{lossless}.



\begin{questionbox}{Domanda 93}
(ModernVC) Why are “Tiles” considered “hardware-friendly” in \textbf{VVC} and \textbf{HEVC}?
\end{questionbox}


Tiles split a frame into independently decodable rectangular regions. They restrict dependencies across boundaries, enabling parallel encoding/decoding on multiple CPU/GPU cores with lower synchronization cost.



\begin{questionbox}{Domanda 94}
(ModernVC) What is the function of an “In-\textbf{loop filter}” like the Adaptive \textbf{loop filter} (\textbf{ALF})?
\end{questionbox}


In-loop filters are applied inside the reconstruction loop before frames are stored as references. This prevents artifacts from being propagated by \textbf{motion compensation}.

H.264 uses deblocking. \textbf{HEVC} adds \textbf{SAO}. \textbf{VVC} adds \textbf{ALF} and other tools.

\textbf{ALF} (Adaptive \textbf{loop filter}) reduces reconstruction artifacts such as ringing and residual \textbf{distortion} by applying adaptive filtering to reconstructed samples before reference storage.



\begin{questionbox}{Domanda 95}
Describe the intra-coding modes in H.264. [Optional] Discuss also the Intra modes in H.265.
\end{questionbox}


\textbf{intra prediction} exploits spatial redundancy from already reconstructed neighboring samples in the same frame.

\begin{itemize}
  \item \textbf{H.264/AVC:} 9 modes for $4\times4$ blocks: 8 directional + DC. For $16\times16$ \textbf{luma} blocks, 4 modes.
  \item \textbf{H.265/\textbf{HEVC}:} 35 modes: 33 directional + DC + Planar.
  \item \textbf{H.266/\textbf{VVC}:} 65 directional modes plus wide-angle variants for non-square blocks.
\end{itemize}

Most Probable Mode (MPM) lists reduce signaling cost by predicting likely modes from neighboring blocks.



\begin{questionbox}{Domanda 96}
Describe the principle of the deblocking in-\textbf{loop filter} of H.264.
\end{questionbox}


At low bitrate, block transform and \textbf{quantization} create visible discontinuities at block boundaries. H.264 applies a normative in-loop \textbf{deblocking filter} on edges between $4\times4$ blocks.

Filtering strength depends on:

\begin{itemize}
  \item coding mode,
  \item motion vectors,
  \item reference frames,
  \item \textbf{quantization} parameter,
  \item local boundary conditions.
\end{itemize}

It is in-loop because filtered reconstructed frames are used as future references; this avoids propagating blocking artifacts.


